\section{Multi-Agent：协作增益与通信成本}

\subsection{何时使用多智能体}
课程并不鼓励 ``为并行而并行''。Multi-agent 适用于任务可分解、角色边界清晰、通信协议稳定的场景。若任务强耦合或中间状态高度共享，单体 Agent + 强 memory 往往更稳。

\begin{knowledgebox}{多智能体收益的来源}
\begin{itemize}
    \item 角色专门化降低单体策略复杂度；
    \item 并行探索提高候选解覆盖率；
    \item 交叉审阅减少单路径幻觉风险。
\end{itemize}
\end{knowledgebox}

\begin{table}[H]
\centering
\begin{tabular}{p{3.3cm}p{3.8cm}p{3.9cm}}
\toprule
\textbf{维度} & \textbf{单体 Agent} & \textbf{Multi-Agent} \\
\midrule
上下文管理 & 简单，但可能拥塞 & 分布式，但需同步协议 \\
执行效率 & 串行，调试容易 & 并行潜力高，协调复杂 \\
错误定位 & 集中式定位较直观 & 需区分角色错误与通信错误 \\
成本结构 & 预测相对稳定 & 通信与重复推理开销高 \\
\bottomrule
\end{tabular}
\caption{单体与多体架构的取舍}
\end{table}

\begin{importantbox}{多智能体系统的落地原则}
先设计协议再增加角色。协议至少包括：消息格式、共享状态范围、冲突解决规则、超时与回退策略。
\end{importantbox}

\begin{warningbox}{多智能体最常见失败模式}
\begin{itemize}
    \item 角色职责重叠，导致重复劳动；
    \item 通信信息膨胀，吞噬上下文预算；
    \item 缺少统一验收器，形成“彼此确认但整体错误”。
\end{itemize}
\end{warningbox}

\subsection{本章小结}
Multi-agent 不是默认更强，而是条件性更强。关键在于任务分解质量和通信协议设计，而非 agent 数量本身。

